\chapter{A first result}

This chapter is a worked example of every blueprint construction you will
need; replace its mathematics with your own. The source is cited as
\cite{Model2020}, and the paper-gap apparatus is demonstrated by
\cite{gap:demo_scope_restriction}.

\begin{theorem}[Squares are nonnegative]
    \label{thm:sq_nonneg}
    \lean{[[ project_name ]].sq_nonneg}
    \leanok
    For every real number \(x\),
    \begin{align}
        0 \le x^2. \label{eq:sq_nonneg}
    \end{align}
\end{theorem}
\begin{proof}\leanok
    The square is a product of two equal factors; a product of two
    nonpositive or two nonnegative factors is nonnegative.
\end{proof}

\begin{theorem}[A dependent result]
    \label{thm:sum_sq_nonneg}
    \lean{[[ project_name ]].sq_nonneg}
    \notready
    \uses{thm:sq_nonneg}
    A sum of squares is nonnegative, by \eqref{eq:sq_nonneg} applied to each
    summand. This entry shows a statement whose proof is not yet
    formalized: it carries \verb|\notready| and its node stays amber in the
    dependency graph until a Lean proof exists.
\end{theorem}
